% chapter1-body.tex — pure content, no preamble.
% This file is \input{} by both main.tex (book) and chapter1.tex (article).




\section{Problem Comprehend}
Consider a one-dimensional line with a wall located at
\[
    x=a>0.
\]
We launch a right-going wave from the left. Its complex representation is
\[
    U(x,t)=e^{ikx}e^{-i\omega t}
    =e^{i(kx-\omega t)}.
\]

Here, $\omega$ is the angular frequency, with unit $\mathrm{s}^{-1}$,
and $k$ is the wavenumber, with unit $\mathrm{m}^{-1}$.
The spatial phase is
\[
    \phi(x)=kx.
\]

For a uniform wave, one complete wavelength $\lambda$ corresponds to
a phase change of $2\pi$. Hence
\[
    k=\frac{2\pi}{\lambda}.
\]

This definition uses a complete wavelength. However, if the phase varies
spatially, a complete wavelength may not be locally available for
measurement. Instead, we consider the phase change over an infinitesimal
spatial interval:
\[
    k(x)
    =\lim_{\Delta x\to0}
    \frac{\phi(x+\Delta x)-\phi(x)}{\Delta x}
    =\frac{d\phi}{dx}.
\]

Thus, the local wavenumber measures the local phase accumulation per
unit distance.

The same idea can also be interpreted through time.
Dividing the local phase change and the propagation distance
by the same time interval gives, intuitively,
\[
    k
    =\frac{d\phi}{dx}
    \sim
    \frac{d\phi/dt}{dx/dt}
    =\frac{\omega}{c}.
\]
Thus, the wavenumber can also be understood as the temporal phase
accumulation rate divided by the propagation speed.

% TODO (Work 2): give a rigorous derivation using the total derivative
% along a constant-phase trajectory.

In higher dimensions, this idea becomes the local wavevector
\[
    \mathbf{k}(x)=\nabla\phi(x).
\]


\subsection{Propagation Directions in One Dimension}
In one dimension, the unit sphere consists of only two points,
\[
    S^0=\{-1,+1\},
\]
corresponding to the two possible propagation directions.

If
\[
    D^1=[-1,1],
\]
then
\[
    \partial D^1=S^0=\{-1,+1\}.
\]

For a direction $d\in S^0$, a plane-wave solution can be written as
\[
    u(x;d)=e^{ikxd}.
\]
Therefore,
\[
    d=+1 \quad\Longrightarrow\quad e^{ikx},
\]
and
\[
    d=-1 \quad\Longrightarrow\quad e^{-ikx}.
\]

% TODO: explain how the discrete directions in $S^0$
% generalize to $d\in S^{n-1}$ in higher dimensions.




\section{Bessel Function Derivation}
Basically I'm going to talk about
derivation of the change of veriable of the Laplacian
\[
    \partial_i=\frac{\partial}{\partial q^i}
    \Delta u=
\]
\[
    c^i=
    \left[\begin{array}{c}
            c^1
            \\
            c^2
        \end{array}\right]^{i}
\]
\[
    p^j=
    \left[\begin{array}{c}
            p^1
            \\
            p^2
        \end{array}\right]^{j}
\]
\[
    u(c^i)=u(p^j)
\]
\[\begin{tikzcd}
        \begin{array}{c}     \left[\begin{array}{c}             c^1             \\             c^2         \end{array}\right]^{i}  \end{array} & u \\
        \begin{array}{c}     \left[\begin{array}{c}             p^1             \\             p^2         \end{array}\right]^{j} \end{array}
        \arrow[from=1-1, to=1-2]
        \arrow["J", from=1-1, to=2-1]
        \arrow["{J^{-1}}", shift left=3, from=2-1, to=1-1]
        \arrow[from=2-1, to=1-2]
    \end{tikzcd}\]
\[
    c^i=\textbf{Polar}(p^j)
\]
Satisfy the Jacobian as follow.
\[
    J=F^i{}_j{}=
    \left[\begin{array}{cc}
            \left[\begin{array}{ll}
                          \frac{\partial c^1}{\partial p^1}
                           &
                          \frac{\partial c^1}{\partial p^2}
                      \end{array}\right]_{j}
            \\
            \left[\begin{array}{ll}
                          \frac{\partial c^2}{\partial p^1}
                           &
                          \frac{\partial c^2}{\partial p^2}
                      \end{array}\right]_{j}
        \end{array}\right]^{i}
    =\frac{\partial c^i}{\partial p^j}
\]
\[
    J^{-1}=B^j{}_i{}=
    \left[\begin{array}{cc}
            \left[\begin{array}{ll}
                          \frac{\partial p^1}{\partial c^1}
                           &
                          \frac{\partial p^1}{\partial c^2}
                      \end{array}\right]_{i}
            \\
            \left[\begin{array}{ll}
                          \frac{\partial p^2}{\partial c^1}
                           &
                          \frac{\partial p^2}{\partial c^2}
                      \end{array}\right]_{i}
        \end{array}\right]^{j}
    =\frac{\partial p^j}{\partial c^i}
\]
\[
    \frac{\partial}{\partial p^j}=
    \left[\begin{array}{ll}
            \frac{\partial }{\partial p^1}
             &
            \frac{\partial }{\partial p^2}
        \end{array}\right]_{j}
    =
    e_j
\]
\[
    \frac{\partial}{\partial c^i}=
    \left[\begin{array}{ll}
            \frac{\partial }{\partial c^1}
             &
            \frac{\partial }{\partial c^2}
        \end{array}\right]_{i}
    =
    \widetilde{e}_i
\]
\[
    \frac{\partial}{\partial p^j}=\frac{\partial}{\partial c^i}F^i{}_j{}
\]
\[
    \frac{\partial}{\partial c^i}=\frac{\partial}{\partial p^j}B^j{}_i{}
\]
then we need to derive the total Gradient of \(F\) and \(B\) as follow
\[
    \frac{\partial}{\partial c^i}B^j{}_i{}=
    \frac{\partial}{\partial c^i}\frac{\partial p^j}{\partial c^i}=
    \left[\begin{array}{cc}
            \frac{\partial^2 p^1}{\partial {c^1}^2}
            +
            \frac{\partial^2 p^1}{\partial {c^2}^2}
            \\
            \frac{\partial^2 p^2}{\partial {c^1}^2}
            +
            \frac{\partial^2 p^2}{\partial {c^2}^2}
        \end{array}\right]^{j}
\]
then Laplacian \(u\) as follow
\[
    \begin{aligned}
        \Delta u                                                    & =\frac{\partial}{\partial c^i}\frac{\partial}{\partial c^i}u                                                                                                                                                     \\
                                                                    & =\frac{\partial}{\partial c^i}\frac{\partial u}{\partial c^i}                                                                                                                                                    \\
                                                                    & =\frac{\partial}{\partial c^i}\left(\frac{\partial p^j}{\partial c^i}\frac{\partial u}{\partial p^j}\right)                                                                                                      \\
                                                                    & =\left(\frac{\partial^2 p^j}{\partial {c^i}^2}\frac{\partial u}{\partial p^j}+\frac{\partial p^j}{\partial c^i}\frac{\partial}{\partial c^i}\frac{\partial u}{\partial p^j}\right)                               \\
        \frac{\partial}{\partial c^i}\frac{\partial}{\partial c^i}u & =\left(\frac{\partial^2 p^j}{\partial {c^i}^2}\frac{\partial u}{\partial p^j}+\frac{\partial p^j}{\partial c^i}\frac{\partial p^k}{\partial c^i}\frac{\partial^2 u}{\partial p^j\partial p^k}\right)             \\
        \frac{\partial}{\partial c^i}\frac{\partial}{\partial c^i}  & =\left(\frac{\partial^2 p^j}{\partial {c^i}^2}\frac{\partial }{\partial p^j}+\frac{\partial p^j}{\partial c^i}\frac{\partial p^k}{\partial c^i}\frac{\partial}{\partial p^j}\frac{\partial}{\partial p^k}\right) \\
    \end{aligned}
\]
\[
    \boxed{
        \Delta=g^{jk}\partial_j\partial_k
        +(\Delta_c p^j)\partial_j
    }
\]
Which
\[
    \boxed{
        A^{jk}=
        \frac{\partial p^j}{\partial c^i}
        \frac{\partial p^k}{\partial c^i}
        =g^{jk}
    }
\]
This \(g^{jk}\) Inverse Metric Tensor
\[
    B^j=\Delta_c p^j
\]
is the Laplacian of coordinate function \(p^j(c)\) under \(c\)




\section{TMP}


\subsection{MVHR}
\[
    T^i{}_{ab}=
    \left[\begin{array}{c}
            \left[\begin{array}{ccc}
                          \left[\begin{array}{NNNN}
                                  1 &
                                  2 &
                                  3 &
                                  4
                              \end{array}\right]_{b}
                           &
                          \left[\begin{array}{NNNN}
                                  5 &
                                  6 &
                                  7 &
                                  8
                              \end{array}\right]_{b}
                           &
                          \left[\begin{array}{NNNN}
                                  9  &
                                  10 &
                                  11 &
                                  12
                              \end{array}\right]_{b}
                      \end{array}\right]_{a}
            \\
            \left[\begin{array}{ccc}
                          \left[\begin{array}{NNNN}
                                  13 &
                                  14 &
                                  15 &
                                  16
                              \end{array}\right]_{b}
                           &
                          \left[\begin{array}{NNNN}
                                  17 &
                                  18 &
                                  19 &
                                  20
                              \end{array}\right]_{b}
                           &
                          \left[\begin{array}{NNNN}
                                  21 &
                                  22 &
                                  23 &
                                  24
                              \end{array}\right]_{b}
                      \end{array}\right]_{a}
        \end{array}\right]_{i}
\]
\[
    T_i{}^a{}_b=
    \left[\begin{array}{cc}
            \left[\begin{array}{ccc}
                          \left[\begin{array}{NNNN}
                                  1 &
                                  2 &
                                  3 &
                                  4
                              \end{array}\right]_{b}
                          \\
                          \left[\begin{array}{NNNN}
                                  5 &
                                  6 &
                                  7 &
                                  8
                              \end{array}\right]_{b}
                          \\
                          \left[\begin{array}{NNNN}
                                  9  &
                                  10 &
                                  11 &
                                  12
                              \end{array}\right]_{b}
                      \end{array}\right]_{a}
             &
            \left[\begin{array}{ccc}
                          \left[\begin{array}{NNNN}
                                  13 &
                                  14 &
                                  15 &
                                  16
                              \end{array}\right]_{b}
                          \\
                          \left[\begin{array}{NNNN}
                                  17 &
                                  18 &
                                  19 &
                                  20
                              \end{array}\right]_{b}
                          \\
                          \left[\begin{array}{NNNN}
                                  21 &
                                  22 &
                                  23 &
                                  24
                              \end{array}\right]_{b}
                      \end{array}\right]_{a}
        \end{array}\right]_{i}
\]
\[
    s=\left[\begin{array}{c}7\end{array}\right]_{0}
\]
\[
    A^i=\left[\begin{array}{c}
            7
            \\
            3
        \end{array}\right]_{i}
\]
\[
    A_i=\left[\begin{array}{cc}7&3\end{array}\right]_{i}
\]
\[
    B^i{}_a=\left[\begin{array}{cc}
            \left[\begin{array}{lll}1&2&3\end{array}\right]_{a}
            \\
            \left[\begin{array}{lll}4&5&6\end{array}\right]_{a}
        \end{array}\right]_{i}
\]
\[
    B_i{}^a=\left[\begin{array}{cc}
            \left[\begin{array}{lll}
                          1
                          \\
                          2
                          \\
                          3\end{array}\right]_{a}
             &
            \left[\begin{array}{lll}
                          1
                          \\
                          2
                          \\
                          3\end{array}\right]_{a}
        \end{array}\right]_{i}
\]
\[
    B^{ia}=\left[\begin{array}{cc}
            \left[\begin{array}{lll}
                          1
                          \\
                          2
                          \\
                          3\end{array}\right]_{a}
            \\
            \left[\begin{array}{lll}
                          1
                          \\
                          2
                          \\
                          3\end{array}\right]_{a}
        \end{array}\right]_{i}
\]
\[
    B_{ia}=\left[\begin{array}{cc}
            \left[\begin{array}{lll}
                          1 &
                          2 &
                          3
                      \end{array}\right]_{a}
             &
            \left[\begin{array}{lll}
                          1 &
                          2 &
                          3
                      \end{array}\right]_{a}
        \end{array}\right]_{i}
\]
TMP




\section{FEM}
三、最重要的通用生成公式
无论几维、几次，我们都从
\[
    (K,P,\Sigma)
\]
开始。

假设
\[
    P=\operatorname{span}\{\ket{p_1},\ldots,\left| p_n \right\rangle\},
\]
是张成向量空间的向量基底

\[
    \Sigma=\{\bra{\ell^1} ,\ldots,\left\langle \ell^n \right|\}.
\]
是泛函基底，也就是对偶空间。

我们希望构造一些自然对偶基底，也就是对偶
\[
    P=\Phi=\operatorname{span}\{\ket{\phi^1},\ldots,\ket{\phi^n}\}
\]
满足
\[
    \braket{\ell^i}{\phi_j}=\delta^{i}{}_{j}.
\]

推导：\\
对任意任意向量\(\ket{p}\)
有
\[
    \ket{p}=\ket{p_j}c^j
\]

作用泛函
\[
    \braket{\ell^i}{p}=\braket{\ell^i}{p_j}c^j=d^i
\]
我们希望
\[
    \ket{p}=\ket{\phi_i}d^i
\]

所以有
\[
    \ket{p}=\ket{\phi_i}d^i=\ket{p_j}c^j
\]
\[
    \braket{\ell^i}{p_j}\ket{\phi_i}=\ket{p_j}
\]
\[
    \ket{\phi_i}=\braket{\ell^i}{p_j}^{-1}\ket{p_j}
\]

构造矩阵：
\[
    \boxed{M^{i}{}_{j}=\braket{\ell^i}{p_j}.}
\]

如果
\[
    \det M\ne0,
\]

就可以定义：
\[
    \boxed{
        \ket{\phi_i}=\sum_{j=1}^n(M^{-1})^{j}{}_{i}\ket{p_j}.
    }
\]

于是直接得到
\[
    \braket{\ell^i}{\phi_j}=\delta^{i}{}_{j}.
\]

这就是所有局部形函数的统一生成器。

有了这组形函数，对任何自由度有定义的函数 \(u\)，就可以建立插值算子：
\[
    \boxed{
        I_Ku=\sum_{j=1}^n\braket{\ell^j}{u}\phi_j.
    }
\]

这一套公式不区分一次元、二次元、Hermite 元、三角形元。

区别只在于我们怎样选择 \(K,P,\Sigma\)。

下面开始具体生成。




\section{Reference}
\bibliographystyle{plain}  % or another style like unsrt, IEEEtran, etc.
\bibliography{refer}  % references.bib is the file name


% \subsection{Definition of EMH}
% \begin{definition}[EMH means Symmetries]
%     \label{def:mydefinition}
%     We say a market is completely EMH if and
%     only if for every condition $A$, the random
%     variable of the logarithmic growth rate $W$
%     agrees that $\mathbb{E}_{W\mid A}(W)\doteq 0$.
% \end{definition}


% \subsection{Prior Work}
% Discussion of what others have done.




% \section{Methods}
% How you approached the problem.




% \section{Results}
% What you found.




% \section{Conclusion}
% Your closing thoughts.




